\begin{center}\textls[120]{d) Die Hundert Theoreme.}\end{center}
\noindent\hspace*{\parindent}Wir wollen hier noch die Besprechung der nachgelassenen Abhandlung
\textit{Hundert Theoreme über die neuen Transzendenten} anschliessen, die auf einzelnen
\par\clearpage\SourceLeaf{267}\ArithmeticHead{182}{L}{L. SCHLESINGER, ÜBER GAUSS' FUNKTIONENTHEORETISCHE ARBEITEN.}
\noindent Blättern (Fi) aufgezeichnet ist, und für deren Entstehungszeit sich als untere
Zeitgrenze 1818 angeben lässt, da darin (S. 467) die \textit{Determinatio attractionis}
zitiert wird; wahrscheinlich ist sie um 1825/26 verfasst. Sie enthält nichts
wesentlich Neues, nur in eleganter und überaus klarer Darstellung Dinge, die
uns schon aus der Scheda Ac und Scheda An bekannt sind.

Für die drei geraden Thetafunktionen werden, da sie »von grosser Wichtigkeit«
sind, »besondere Funktionalzeichen« \(P(x,y),\ Q(x,y),\ R(x,y)\) eingeführt
(vergl. oben, S. 76), und in den drei ersten Theoremen Identitäten zwischen
Reihen und Produkten hergeleitet. Es folgt (vergl. oben S. 82, Fussnote \textsuperscript{110)})
als viertes Theorem die Identität (3) Abschn. II (oben S. 16), die die Umformung
der Reihe für \(P(x,y)\) in ein unendliches Produkt liefert, und als
fünftes die analoge Identität, die dasselbe für die Funktion \(R(x,y)\) leistet.
Für \(y=1\) ergeben sich die Reihen \(p,q,r\); über die für diese aufgestellten
Sätze (die Zählung der Theoreme hört hier auf) ist bereits oben Abschn. V c)
berichtet worden.

Auf dieselbe Zeit wie die \textit{Hundert Theoreme} würde die Notiz zu datieren
sein, die Werke III, S. 384, abgedruckt ist und aus dem Handbuch 19, Be
(begonnen November 1809), S. 147 stammt. Es finden sich nämlich (vergl.
die Angabe von \textsc{Schering} a. a. O.) auf der vorhergehenden Seite (146) Formeln
über die Reihen \(P(x,y),\ Q(x,y),\ R(x,y)\), also mit denselben Bezeichnungen
wie in den \textit{Hundert Theoremen} aufgezeichnet. Da auf S. 136 des Handbuchs
das Datum »1825, Dec. 4.« steht, ist die in Rede stehende Notiz sicher nach
diesem Tage geschrieben. Sie enthält das Theorem: »das Produkt zweier
Summen von vier Quadraten ist selbst eine Summe von vier Quadraten«, aus
dem bekanntlich \textsc{Jacobi} in seinen Vorlesungen die wichtigsten Eigenschaften
der Thetareihen abgeleitet hat\footnote{Siehe \textsc{Jacobi}s Werke I, S. 503 ff. Ob \textsc{Gauss} wirklich, wie \textsc{Schering} vermutet, mit Hilfe dieses Theorems die Darstellbarkeit von zwei beliebigen Zahlen \(m,n\) durch die Funktionen \(p(x)^2,q(x)^2\) für
\[
\log x=-\pi\frac{M(m,n)}{M(m,\sqrt{m^2-n^2})}
\]
nachgewiesen hat, entzieht sich unserer Kenntnis, jedenfalls wissen wir, dass er diese Darstellbarkeit auf andere Weise gefunden hat (siehe Abschn. V, c), d)).}.

In den Jahren 1825—1826 scheinen sich die Verhältnisse in Göttingen
auch wieder mehr nach den berechtigten Wünschen von \textsc{Gauss} gestaltet zu
\par\clearpage\SourceLeaf{268}\ArithmeticHead{183}{R}{IXa. ABELS UND JACOBIS ENTDECKUNGEN.}
\noindent haben. Die sich durch mehrere Jahre hinziehenden Verhandlungen über eine
Berufung nach Berlin wurden, weil einerseits die dort einzugehenden Verpflichtungen
allzu drückend erschienen, und andererseits die hannöversche Regierung
auch bedeutendes Entgegenkommen zeigte, von \textsc{Gauss}' Seite abgebrochen;
»ich bin jetzt«, schreibt er am 21. März 1825 an \textsc{Pfaff}\footnote{\textit{Sammlung von Briefen \mbox{.\,.} zwischen J. Fr. Pfaff und \mbox{.\,.} andern}, Leipzig 1853, S. 278, zitiert bei \textsc{Sartorius}, a. a. O. \textsuperscript{13)}, S. 60.}, »Konjunkturen
abgerechnet, die ausser aller Berechnung liegen, fürs Leben an
Göttingen geknüpft \mbox{.\,.\,.} durch das Band aufrichtiger Dankbarkeit für sehr
liberales Benehmen unserer Regierung«. Und im Jahre 1826 fanden auch die
praktisch-geodätischen Arbeiten im Gelände ihren Abschluss, so dass in den
nächstfolgenden Jahren wieder mehr Zeit und Kraft auf die Ausführung theoretischer
Untersuchungen verwendet werden konnte.
